\par\Needspace{6\baselineskip}\medskip
\begin{center}\hypertarget{native-168-article121}{}\pdfbookmark[2]{121.}{article-121}121.\end{center}
Liegen bereits Näherungswerte der Unbekannten vor, so lassen sich daraus nach dem eben erläuterten Verfahren die wahren Werte mit aller gewünschten Genauigkeit \SourceLeaf{169}herleiten. Zuerst berechnen wir die Werte von $X,Y$, die diesen Näherungswerten $a,b$ entsprechen. Verschwinden sie dabei nicht schon von selbst, so wiederholen wir die Rechnung mit einem anderen Wertepaar $a',b'$, das vom ersten nur wenig abweicht, und danach mit einem dritten Wertepaar $a'',b''$, sofern nicht $X$ und $Y$ beim zweiten Paar zufällig verschwunden sind. Dann erhalten wir nach den Formeln des vorhergehenden Artikels die wahren Werte, soweit die Annahme, auf der diese Formeln beruhen, von der Wahrheit nicht merklich abweicht. Um dies besser beurteilen zu können, berechnen wir die Werte von $X,Y$ nochmals mit den berichtigten Werten. Zeigt diese Rechnung, dass die Gleichungen $X=0$, $Y=0$ noch nicht erfüllt sind, so ergeben sich doch gewiss viel kleinere Werte von $X,Y$ als aus den ersten drei Hypothesen. Die so erhaltenen Bahnelemente sind daher weit genauer als die den ersten Hypothesen entsprechenden. Wollen wir weitergehen, so lassen wir am besten die Hypothese weg, die die größten Differenzen ergeben hat, und verbinden die beiden übrigen mit der vierten. So bilden wir nach dem vorhergehenden Artikel ein fünftes Wertepaar von $x,y$. Ebenso können wir, wo es sich der Mühe zu lohnen scheint, zu einer sechsten Hypothese und so weiter übergehen, bis die Gleichungen $X=0$, $Y=0$ so genau erfüllt sind, wie es die logarithmischen und trigonometrischen Tafeln erlauben. Nur sehr selten wird es jedoch nötig sein, über das vierte Wertepaar hinauszugehen, sofern die ersten Hypothesen nicht noch zu weit von der Wahrheit entfernt waren.
